\documentclass[12pt, a4paper]{article}

%% preamble.tex — Package loading, custom commands, and style settings
%% All settings that belong before \begin{document} live here.

%% ── Encoding & fonts ──────────────────────────────────────────────────────────
\usepackage[T1]{fontenc}
\usepackage[utf8]{inputenc}
\usepackage{microtype}          % Improved text justification and kerning

%% ── Page geometry ─────────────────────────────────────────────────────────────
\usepackage[
	a4paper,
	top=25mm, bottom=25mm,
	left=30mm, right=25mm
]{geometry}

%% ── Mathematics ───────────────────────────────────────────────────────────────
\usepackage{amsmath}
\usepackage{amssymb}

%% ── Graphics ──────────────────────────────────────────────────────────────────
\usepackage{graphicx}
\graphicspath{{figures/}}

%% ── Tables ────────────────────────────────────────────────────────────────────
\usepackage{booktabs}           % Professional-quality horizontal rules
\usepackage{array}
\usepackage{tabularx}
\usepackage{makecell}           % Multi-line header cells via \makecell[r]{...}

%% ── Colors ────────────────────────────────────────────────────────────────────
\usepackage[dvipsnames, table]{xcolor}
%% styles/my_style.sty — Project-specific style definitions
%% Loaded via %% styles/my_style.sty — Project-specific style definitions
%% Loaded via \input{styles/my_style} from preamble.tex.
%% Do NOT add \usepackage calls here; use preamble.tex instead.

%% ── Project color palette ─────────────────────────────────────────────────────
\definecolor{OceanDeep}{HTML}{003F72}
\definecolor{OceanMid}{HTML}{0077B6}
\definecolor{OceanLight}{HTML}{ADE8F4}
\definecolor{CoralAccent}{HTML}{E07A5F}
\definecolor{NoteGray}{HTML}{6B7280}

%% ── Theorem-style environments ────────────────────────────────────────────────
\usepackage{amsthm}

\newtheoremstyle{definition-style}
  {6pt}{6pt}{\normalfont}{}
  {\bfseries\color{OceanDeep}}{.}{0.5em}{}

\theoremstyle{definition-style}
\newtheorem{definition}{Definition}[section]
\newtheorem{hypothesis}{Hypothesis}[section]

\newtheoremstyle{remark-style}
  {4pt}{4pt}{\normalfont}{}
  {\itshape\color{NoteGray}}{.}{0.5em}{}

\theoremstyle{remark-style}
\newtheorem*{remark}{Remark}

%% ── Section heading style ─────────────────────────────────────────────────────
\usepackage{titlesec}

\titleformat{\section}
  {\large\bfseries\color{OceanDeep}}
  {\thesection}{1em}{}

\titleformat{\subsection}
  {\normalsize\bfseries\color{OceanMid}}
  {\thesubsection}{1em}{}
 from preamble.tex.
%% Do NOT add \usepackage calls here; use preamble.tex instead.

%% ── Project color palette ─────────────────────────────────────────────────────
\definecolor{OceanDeep}{HTML}{003F72}
\definecolor{OceanMid}{HTML}{0077B6}
\definecolor{OceanLight}{HTML}{ADE8F4}
\definecolor{CoralAccent}{HTML}{E07A5F}
\definecolor{NoteGray}{HTML}{6B7280}

%% ── Theorem-style environments ────────────────────────────────────────────────
\usepackage{amsthm}

\newtheoremstyle{definition-style}
  {6pt}{6pt}{\normalfont}{}
  {\bfseries\color{OceanDeep}}{.}{0.5em}{}

\theoremstyle{definition-style}
\newtheorem{definition}{Definition}[section]
\newtheorem{hypothesis}{Hypothesis}[section]

\newtheoremstyle{remark-style}
  {4pt}{4pt}{\normalfont}{}
  {\itshape\color{NoteGray}}{.}{0.5em}{}

\theoremstyle{remark-style}
\newtheorem*{remark}{Remark}

%% ── Section heading style ─────────────────────────────────────────────────────
\usepackage{titlesec}

\titleformat{\section}
  {\large\bfseries\color{OceanDeep}}
  {\thesection}{1em}{}

\titleformat{\subsection}
  {\normalsize\bfseries\color{OceanMid}}
  {\thesubsection}{1em}{}
    % Project-specific color palette and environments

%% ── Bibliography (biber backend) ──────────────────────────────────────────────
\usepackage[
	backend=biber,
	style=authoryear,
	sorting=nyt,
	natbib=true
]{biblatex}
\addbibresource{references.bib}

%% ── Cross-references ──────────────────────────────────────────────────────────
\usepackage[hidelinks]{hyperref}
\usepackage[capitalise, noabbrev]{cleveref}

%% ── Utilities ─────────────────────────────────────────────────────────────────
\usepackage{csquotes}           % Recommended by biblatex
\usepackage{lipsum}             % Placeholder text (remove in final version)
\usepackage{setspace}
\onehalfspacing

%% ── Custom commands ───────────────────────────────────────────────────────────

% Species name: italic, properly typeset
\newcommand{\spname}[1]{\textit{#1}}

% Frequency unit shorthand
\newcommand{\kHz}{\,\text{kHz}}
\newcommand{\Hz}{\,\text{Hz}}

% Shorthand for confidence interval
\newcommand{\ci}[2]{[#1,\,#2]}

% Sequence notation: acoustic event label in small caps
\newcommand{\aseq}[1]{\textsc{#1}}


\begin{document}

\title{Pre-Departure Acoustic Farewell Sequences in
	\emph{Tursiops truncatus}:\\
	A Computational Analysis}
\author{%
	Arthur Dent\thanks{Department of Applied Xenolinguistics, Betelgeuse
		University, Earth (Third Planet), Sol System.}
	\and
	Ford Prefect\thanks{Hitchhiker's Guide Research Division, Magrathea.}
}
\date{May 2026}

\maketitle

\begin{abstract}
	Despite decades of research into cetacean communication, the acoustic
	structure of departure-specific vocalizations in bottlenose dolphins
	(\spname{Tursiops truncatus}) has received surprisingly little systematic
	attention. This paper presents a computational linguistic analysis of
	42~recorded pre-departure sequences from three Atlantic populations,
	demonstrating statistically significant structural differences between
	farewell and non-farewell vocalizations. Our findings suggest that
	bottlenose dolphins possess a dedicated acoustic repertoire for
	leave-taking events, potentially encoding contextual information about
	destination and urgency. We discuss implications for interspecies
	communication theory and the broader question of linguistic intentionality
	in non-human species.
\end{abstract}

\newpage
\tableofcontents
\newpage

% Read STYLE.md before editing prose in this file.
% Follow the Content Creation Workflow in WORKFLOW.md §Content Creation Workflow.

\section{Introduction}
\label{sec:intro}

Dolphins have long been recognized as among the most communicatively
sophisticated non-human animals on Earth~\citep{Wonko2021communication}.
Their vocalizations span a wide frequency range---from low-frequency
burst-pulse sounds near 200\Hz{} to high-frequency whistles exceeding
20\kHz{}---and serve functions including individual identification,
group coordination, and echolocation~\citep{Trillian2023signal}.
Yet despite this richness, one category of dolphin vocalization has
attracted remarkably little systematic study: the acoustic signals
produced immediately before a group or individual permanently departs a
location.

\begin{hypothesis}[Farewell Specificity]
  \label{hyp:specificity}
  Pre-departure vocalizations in \spname{Tursiops truncatus} form a
  structurally distinct class, separable from non-departure vocalizations
  by measurable acoustic features including frequency contour,
  inter-pulse interval, and burst duration.
\end{hypothesis}

This paper tests \cref{hyp:specificity} by applying computational
linguistics methods---specifically hidden Markov model (HMM) sequence
alignment and spectral envelope analysis---to a curated corpus of
42~departure events recorded across three Atlantic coastal populations
between 2021 and 2024~\citep{GuideResearch2025dataset}.

\subsection{Contributions}
\label{subsec:contrib}

The main contributions of this work are as follows.

\begin{itemize}
  \item We introduce the \emph{Atlantic Cetacean Farewell Corpus}
        (ACFC v1.0), the first annotated dataset of departure-specific
        \spname{T.\ truncatus} vocalizations.

  \item We identify three recurring sequence archetypes---which we term
        \aseq{brief-farewell}, \aseq{extended-farewell}, and
        \aseq{urgent-departure}---and characterize their acoustic profiles
        (see \cref{tab:vocalization-comparison}).

  \item We demonstrate that a bigram HMM trained on ACFC achieves
        86.4\% accuracy in distinguishing departure from non-departure
        sequences, substantially above the chance baseline of 50\%.
\end{itemize}

\subsection{Paper Organization}
\label{subsec:organization}

\Cref{sec:background} reviews related literature on cetacean
communication and computational approaches to bioacoustic classification.
\Cref{sec:methodology} describes the data collection protocol,
annotation scheme, and computational methods.
We conclude with a discussion of implications for interspecies
communication theory.

% Read STYLE.md before editing prose in this file.
% Follow the Content Creation Workflow in WORKFLOW.md §Content Creation Workflow.

\section{Background}
\label{sec:background}

\subsection{Cetacean Acoustic Communication}
\label{subsec:cetacean-acoustics}

Bottlenose dolphins (\spname{Tursiops truncatus}) produce three broad
classes of sound: (1)~frequency-modulated whistles, which serve as
individual identification signals or ``signature whistles'';
(2)~burst-pulse sounds, a rapid series of clicks used in social
contexts; and (3)~click trains for echolocation~\citep{Wonko2021communication}.
Research into the semantic content of these vocalizations has accelerated
over the past two decades, yet leave-taking behavior has been studied
primarily from an ethological rather than an acoustic perspective.

\citet{Slartibartfast2023topology} note, almost in passing, that coastal
populations near fjord environments exhibit unusually structured departure
calls, attributing this to the particular resonant properties of such
coastlines.
This observation motivated our own data collection strategy, described
in \cref{sec:methodology}.

\subsection{Computational Bioacoustics}
\label{subsec:comp-bioacoustics}

Sequence-level analysis of animal vocalizations has benefited greatly
from methods developed for human speech processing.
Hidden Markov models~(HMMs) were first applied to bird song classification
by \citet{Deep2022number}, who demonstrated that temporal dependencies
within a vocalization bout could be captured by the Markov transition
matrix, achieving state-of-the-art accuracy on a 12-species corpus.

More recently, \citet{Trillian2023signal} applied Fourier-based spectral
envelope analysis to burst-pulse sequences, showing that departure-adjacent
bouts in captive dolphins exhibited measurably broader spectral envelopes
than those recorded during feeding events.
Our work extends this result to wild populations and tests whether the
pattern generalizes across geographically distinct groups.

\begin{figure}[htbp]
	\centering
	% TODO: Replace the placeholder below with the actual spectrogram image:
	%   figures/ch01/spectrogram_example.pdf
	\fbox{\rule{0pt}{4cm}\hspace{8cm}}%
	\caption{Representative spectrogram of an \aseq{extended-farewell}
		sequence recorded at Site~A (Atlantic, 2023). The vertical axis
		shows frequency in \kHz{}; the horizontal axis shows time in
		seconds. Color indicates power spectral density.
		(Placeholder: replace with \texttt{figures/ch01/spectrogram\_example.pdf}.)
	}
	\label{fig:ch01-spectrogram-example}
\end{figure}

\Cref{fig:ch01-spectrogram-example} illustrates a representative
\aseq{extended-farewell} sequence; note the characteristic frequency
sweep from approximately 8\kHz{} to 18\kHz{} occurring in the
final 1.2~seconds before departure.

\subsection{Leave-Taking in Animal Communication}
\label{subsec:leave-taking}

The study of leave-taking rituals in animal communication is a nascent
field~\citep{Marvin2022depression}.
Across species ranging from great apes to corvids, researchers have
documented the existence of structured terminal vocalizations that differ
acoustically from vocalizations produced in non-departure contexts.
\citet{Prefect2025syntax} provides the most direct antecedent to our
work: a minimalist syntactic analysis of 17~\spname{Tursiops} farewell
calls recorded near the Azores, arguing for a recursive hierarchical
structure in the calls' temporal organization.
Our study corroborates Prefect's syntactic findings while providing the
first corpus-scale statistical evidence for their generality.

% Read STYLE.md before editing prose in this file.
% Follow the Content Creation Workflow in WORKFLOW.md §Content Creation Workflow.

\section{Methodology}
\label{sec:methodology}

\subsection{Data Collection}
\label{subsec:data-collection}

Acoustic recordings were obtained at three Atlantic sites: Site~A
(western Azores archipelago), Site~B (Bay of Biscay, northern Spain),
and Site~C (Cape Verde islands).
Field observations were conducted between June~2021 and September~2024,
yielding 42~confirmed departure events for which full acoustic records
were available.
A departure event was defined as a session in which a marked individual
or coherent group left the observation area and did not return within
the 72-hour monitoring window.

\begin{definition}[Departure Event]
	\label{def:departure-event}
	A \emph{departure event} is a continuous acoustic recording window of
	\(\pm 5\) minutes around the last observed surfacing of a focal
	individual or group, provided that individual or group did not re-enter
	the monitoring area within 72~hours of the recording.
\end{definition}

Recordings were made using calibrated hydrophone arrays deployed at
depths of 3--8~m, sampling at 96\kHz{} with 24-bit resolution.
All recordings were stored in lossless FLAC format and are archived
with the \citet{GuideResearch2025dataset} dataset.

\subsection{Annotation Scheme}
\label{subsec:annotation}

Each departure event was segmented into discrete acoustic units by two
independent annotators following the protocol in
\citet{GuideResearch2025dataset}.
Inter-annotator agreement, measured by Cohen's~$\kappa$, was $0.83$
across the full corpus---within the range typically considered
``substantial'' agreement~\citep{Dent2024farewell}.

Each unit was assigned one of five labels:
\texttt{whistle}, \texttt{burst-pulse}, \texttt{click-train},
\texttt{silence}, or \texttt{noise}.
The sequence of labels for a departure event constitutes a
\emph{vocalization trace}.

\subsection{Computational Analysis}
\label{subsec:analysis}

\paragraph{Hidden Markov Model}

We model each vocalization trace as a sequence emitted by a discrete-time
HMM with $N$ hidden states representing acoustic context:

\begin{equation}
	\label{eq:hmm-likelihood}
	P(\mathbf{o} \mid \lambda) =
	\sum_{\mathbf{q}} P(\mathbf{o} \mid \mathbf{q}, \lambda)\,
	P(\mathbf{q} \mid \lambda),
\end{equation}

where $\mathbf{o} = (o_1, \ldots, o_T)$ is the observed label sequence,
$\mathbf{q} = (q_1, \ldots, q_T)$ is the hidden state sequence, and
$\lambda = (A, B, \pi)$ denotes the transition matrix, emission matrix,
and initial distribution, respectively.

Model parameters were estimated using the Baum--Welch algorithm with
random restarts ($k=50$).
The number of hidden states $N$ was selected by Bayesian information
criterion over the range $N \in \{2, 3, \ldots, 10\}$; the
optimal value was $N = 4$.

\paragraph{Spectral Envelope Analysis}

For each annotated whistle unit, we extracted the spectral envelope using
linear predictive coding (LPC) with order $p = 16$.
The resulting envelope coefficients were projected onto a 2D space via
principal component analysis (PCA) for visualization.

\Cref{tab:vocalization-comparison} summarizes the acoustic features of
the three identified archetypes.

%% tables/comparison_table.tex — Vocalization archetype comparison table
%% Included via %% tables/comparison_table.tex — Vocalization archetype comparison table
%% Included via \input{tables/comparison_table} in chapters/03_methodology.tex

\begin{table}[htbp]
  \caption{Acoustic features of the three departure vocalization archetypes
    identified in the Atlantic Cetacean Farewell Corpus (ACFC v1.0).
    Values are mean $\pm$ standard deviation across all instances.
    Frequency ranges are for the dominant whistle component.}
  \label{tab:vocalization-comparison}
  \begin{tabularx}{\linewidth}{>{{\raggedright\arraybackslash}}X r r r >{\raggedright\arraybackslash}X}
    \toprule
    \textbf{Archetype}
      & \makecell[r]{\textbf{Peak Freq.}\\\textbf{(\kHz)}}
      & \textbf{Duration (s)}
      & \textbf{Count}
      & \textbf{Characteristic Feature} \\
    \midrule
    \aseq{brief-farewell}
      & $14.2 \pm 1.8$
      & $0.83 \pm 0.21$
      & 18
      & Single upward frequency sweep; minimal burst-pulse content \\[3pt]
    \aseq{extended-farewell}
      & $11.6 \pm 2.3$
      & $3.47 \pm 0.94$
      & 17
      & Multi-phrase structure with recurring motif; broad spectral envelope \\[3pt]
    \aseq{urgent-departure}
      & $17.1 \pm 0.9$
      & $0.41 \pm 0.11$
      & 7
      & High-frequency, short-duration; high burst-pulse density \\
    \bottomrule
  \end{tabularx}
\end{table}
 in chapters/03_methodology.tex

\begin{table}[htbp]
  \caption{Acoustic features of the three departure vocalization archetypes
    identified in the Atlantic Cetacean Farewell Corpus (ACFC v1.0).
    Values are mean $\pm$ standard deviation across all instances.
    Frequency ranges are for the dominant whistle component.}
  \label{tab:vocalization-comparison}
  \begin{tabularx}{\linewidth}{>{{\raggedright\arraybackslash}}X r r r >{\raggedright\arraybackslash}X}
    \toprule
    \textbf{Archetype}
      & \makecell[r]{\textbf{Peak Freq.}\\\textbf{(\kHz)}}
      & \textbf{Duration (s)}
      & \textbf{Count}
      & \textbf{Characteristic Feature} \\
    \midrule
    \aseq{brief-farewell}
      & $14.2 \pm 1.8$
      & $0.83 \pm 0.21$
      & 18
      & Single upward frequency sweep; minimal burst-pulse content \\[3pt]
    \aseq{extended-farewell}
      & $11.6 \pm 2.3$
      & $3.47 \pm 0.94$
      & 17
      & Multi-phrase structure with recurring motif; broad spectral envelope \\[3pt]
    \aseq{urgent-departure}
      & $17.1 \pm 0.9$
      & $0.41 \pm 0.11$
      & 7
      & High-frequency, short-duration; high burst-pulse density \\
    \bottomrule
  \end{tabularx}
\end{table}



\printbibliography[heading=bibintoc, title={References}]

\end{document}
